\documentclass{article}
\usepackage[ngerman]{babel}
\usepackage[T1]{fontenc}
\usepackage[utf8]{inputenc}
\usepackage{amsmath}

\title{Zur Approximation ganzer Funktionen durch Polynome}
\author{K.~M\"uller}
\date{2026}

\begin{document}
\maketitle

\section{Einleitung}

Die Approximationstheorie besch\"aftigt sich mit der Frage, wie gut sich
komplizierte Funktionen durch einfachere ann\"ahern lassen.  In dieser Arbeit
untersuchen wir die gleichm\"a\ss{}ige Konvergenz von Polynomfolgen auf
kompakten Teilmengen der komplexen Ebene.

% LANG-012 / DE-006: wrong quote style (English quotes instead of German)
Die sogenannte "Weierstra\ss-Approximation" ist ein klassisches Ergebnis.
Man spricht von "gleichm\"a\ss{}iger Konvergenz" im Banachraum.

% Correct: German quotes
Die sogenannte \glqq Weierstra\ss-Approximation\grqq{} bildet die Grundlage.
Man spricht von \glqq gleichm\"a\ss{}iger Konvergenz\grqq{} im Funktionenraum.

\section{Grundlagen}

Sei $f \colon \mathbb{C} \to \mathbb{C}$ eine ganze Funktion.  Nach dem Satz
von Liouville ist jede beschr\"ankte ganze Funktion konstant.  D.\,h.\ wenn
$|f(z)| \le M$ f\"ur alle $z \in \mathbb{C}$, dann ist $f$ konstant.

% German abbreviations
Bzgl.\ der Konvergenz unterscheidet man verschiedene Typen.  Z.\,B.\ ist die
punktweise Konvergenz schw\"acher als die gleichm\"a\ss{}ige.  U.\,a.\ gilt
der Satz von Dini f\"ur monotone Folgen.  S.\,o.\ f\"ur weitere Details.

% Correct use of \ss
Die Gr\"o\ss{}e des Fehlers h\"angt von der Gl\"atte der Funktion ab.  Au\ss{}erdem
spielt die Gestalt des Definitionsbereichs eine wesentliche Rolle.  Das Ma\ss{}
der Abweichung wird in der Supremumsnorm gemessen.

\section{Hauptergebnis}

\begin{theorem}
F\"ur jede stetige Funktion $f \colon [a,b] \to \mathbb{R}$ und jedes
$\varepsilon > 0$ existiert ein Polynom $p$ mit
\[
  \|f - p\|_\infty < \varepsilon.
\]
\end{theorem}

Der Beweis verl\"auft \"uber die Bernstein-Polynome
\[
  B_n(f; x) = \sum_{k=0}^{n} \binom{n}{k} f\!\left(\frac{k}{n}\right)
  x^k (1-x)^{n-k}.
\]

Die Konvergenz $B_n(f) \to f$ ist gleichm\"a\ss{}ig auf $[0,1]$.

\section{Zusammenfassung}

Wir haben gezeigt, dass die Bernstein-Polynome ein konstruktives Werkzeug
zur gleichm\"a\ss{}igen Approximation stetiger Funktionen darstellen.
Weiterf\"uhrende Untersuchungen betreffen die Konvergenzgeschwindigkeit
und die Verallgemeinerung auf h\"oherdimensionale Gebiete.

\end{document}
